\documentclass[10pt,a4paper]{article}
\usepackage[utf8]{inputenc}
\usepackage[vietnamese]{babel}
\usepackage{amsmath,amssymb,amsfonts}
\usepackage{geometry}
\geometry{a4paper, left=10mm, right=10mm, top=10mm, bottom=10mm}
\usepackage{tikz}
\usetikzlibrary{shapes,arrows.meta,positioning,calc}
\usepackage[most]{tcolorbox}
\usepackage{array,booktabs,multicol}
\usepackage{fontawesome}
\usepackage{helvet}
\renewcommand{\familydefault}{\sfdefault}
\usepackage{parskip}

\definecolor{primary}{RGB}{112, 48, 160}
\definecolor{teal}{RGB}{0, 128, 128}
\definecolor{orange}{RGB}{230, 108, 44}
\definecolor{green}{RGB}{46, 139, 87}
\definecolor{rose}{RGB}{192, 0, 0}

\pagestyle{empty}

\begin{document}

% --- HEADER BANNER ---
\begin{tcolorbox}[
  enhanced,
  colback=violet!8!white,
  colframe=violet!80!black,
  arc=2.5mm, boxrule=1pt,
  left=3mm, right=3mm, top=2mm, bottom=2mm
]
  \begin{tabular*}{\textwidth}{@{\extracolsep{\fill}} lr}
    {\small\bfseries\color{violet!85!black}\faBook\ SỔ TAY KIẾN THỨC TOÁN 2} & 
    {\small\bfseries\color{teal!85!black}\faGraduationCap\ CHƯƠNG TRÌNH GDPT 2018}
  \end{tabular*}
  \vspace{-1mm}
  \begin{center}
    \rule{0.98\textwidth}{0.6pt}
  \end{center}
  \vspace{-2mm}
  \centering
  {\LARGE\bfseries\color{violet!90!black}\faBookmark\ BÀI 71: ÔN TẬP PHÉP NHÂN, PHÉP CHIA}\\[1mm]
  {\itshape\footnotesize\color{black!75} Tổng ôn bảng nhân 2, bảng nhân 5, bảng chia 2, bảng chia 5 và thành phần phép tính}
\end{tcolorbox}

\vspace{-2mm}

\noindent
\begin{minipage}[t]{0.49\textwidth}\vspace{0pt}
\begin{tcolorbox}[
  enhanced,
  colback=violet!4!white,
  colframe=violet!75!black,
  arc=2mm, boxrule=0.8pt,
  left=2mm, right=2mm, top=1.8mm, bottom=1.8mm,
  title={\footnotesize\bfseries\color{white}\faKey\ 1. Bảng nhân 2, 5 và Bảng chia 2, 5}
]
{\footnotesize
\begin{itemize}\setlength{\itemsep}{1pt}\setlength{\leftskip}{-2.5mm}
  \item \textbf{Bảng nhân 2:} Tích là các số chẵn tăng dần $2$ đơn vị:
  $$2, 4, 6, 8, 10, 12, 14, 16, 18, 20$$
  \item \textbf{Bảng nhân 5:} Tích có tận cùng là $0$ hoặc $5$:
  $$5, 10, 15, 20, 25, 30, 35, 40, 45, 50$$
  \item \textbf{Quan hệ nhân $\leftrightarrow$ chia:} Từ một phép nhân ta viết được hai phép chia tương ứng:
  $$a \times b = c \implies c : a = b \text{ và } c : b = a$$
  \textit{Ví dụ:} $2 \times 7 = 14 \implies 14 : 2 = 7$ và $14 : 7 = 2$.
\end{itemize}
}
\end{tcolorbox}
\end{minipage}
\hfill
\begin{minipage}[t]{0.49\textwidth}\vspace{0pt}
\begin{tcolorbox}[
  enhanced,
  colback=teal!4!white,
  colframe=teal!75!black,
  arc=2mm, boxrule=0.8pt,
  left=2mm, right=2mm, top=1.8mm, bottom=1.8mm,
  title={\footnotesize\bfseries\color{white}\faSliders\ 2. Tên gọi thành phần phép tính}
]
{\footnotesize
\begin{itemize}\setlength{\itemsep}{1pt}\setlength{\leftskip}{-2.5mm}
  \item \textbf{Phép nhân:}
  $$\underbrace{a}_{\text{Thừa số}} \times \underbrace{b}_{\text{Thừa số}} = \underbrace{c}_{\text{Tích}}$$
  \item \textbf{Phép chia:}
  $$\underbrace{a}_{\text{Số bị chia}} : \underbrace{b}_{\text{Số chia}} = \underbrace{c}_{\text{Thương}}$$
  \item \textbf{Lưu ý đặc biệt:}
  \begin{itemize}\setlength{\itemsep}{0.5pt}\setlength{\leftskip}{-2mm}
    \item Số nào nhân với $1$ cũng bằng chính số đó.
    \item Số nào chia cho $1$ cũng bằng chính số đó.
    \item Số $0$ nhân hoặc chia cho số nào khác $0$ cũng bằng $0$.
    \item \textbf{Không có phép chia cho 0!}
  \end{itemize}
\end{itemize}
}
\end{tcolorbox}
\end{minipage}

\vspace{-1.5mm}
\begin{tcolorbox}[
  enhanced,
  colback=orange!4!white,
  colframe=orange!80!black,
  arc=2mm, boxrule=0.8pt,
  left=2mm, right=2mm, top=1.5mm, bottom=1.5mm,
  title={\footnotesize\bfseries\color{white}\faSitemap\ 3. Sơ đồ liên kết nhân đôi và chia đều}
]
\centering
\begin{tikzpicture}[x=1.55cm, y=0.75cm]
  \tikzset{arr/.style={-{Stealth[scale=0.8]}}, every node/.append style={align=center}}
  \node[rectangle, rounded corners=2pt, draw=violet!80!black, fill=violet!15, font=\scriptsize\bfseries] (m) at (0,0) {BẢNG NHÂN 2 \& 5\\Gấp lên $2$ lần, $5$ lần};
  \node[rectangle, rounded corners=2pt, draw=teal!80!black, fill=teal!15, font=\scriptsize\bfseries] (d) at (3.5,0) {BẢNG CHIA 2 \& 5\\Chia thành $2$ phần, $5$ phần};

  \draw[arr, thick, bend left=25] (m) to node[above, font=\tiny\bfseries] {Phép tính ngược} (d);
  \draw[arr, thick, bend left=25] (d) to node[below, font=\tiny\bfseries] {Thử lại bằng tích} (m);
\end{tikzpicture}
\end{tcolorbox}

\vspace{-1.5mm}
\noindent
\begin{minipage}[t]{0.49\textwidth}\vspace{0pt}
\begin{tcolorbox}[
  enhanced,
  colback=green!4!white,
  colframe=green!60!black,
  arc=2mm, boxrule=0.8pt,
  left=2mm, right=2mm, top=1.8mm, bottom=1.8mm,
  title={\footnotesize\bfseries\color{white}\faCheckCircle\ 4. Bài tập thực hành mẫu}
]
{\footnotesize
\textbf{1. Tính nhẩm nhanh:}\\
- $2 \times 8 = \mathbf{16}$; \quad $18 : 2 = \mathbf{9}$.\\
- $5 \times 7 = \mathbf{35}$; \quad $45 : 5 = \mathbf{9}$.\\
- $2 \times 4 \times 5 = 8 \times 5 = \mathbf{40}$.\\
\textbf{2. Bài toán có lời văn:}\\
Có $35$ quyển vở chia đều cho $5$ bạn học sinh giỏi. Hỏi mỗi bạn nhận được bao nhiêu quyển vở?\\
\textit{Giải:} Mỗi bạn nhận được số quyển vở là:\\
$35 : 5 = \mathbf{7}$ (quyển vở).
}
\end{tcolorbox}
\end{minipage}
\hfill
\begin{minipage}[t]{0.49\textwidth}\vspace{0pt}
\begin{tcolorbox}[
  enhanced,
  colback=rose!4!white,
  colframe=rose!70!black,
  arc=2mm, boxrule=0.8pt,
  left=2mm, right=2mm, top=1.8mm, bottom=1.8mm,
  title={\footnotesize\bfseries\color{white}\faExclamationTriangle\ 5. Lỗi phân biệt dạng toán}
]
{\footnotesize
\begin{itemize}\setlength{\itemsep}{1.5pt}\setlength{\leftskip}{-2.5mm}
  \item \textbf{Phân biệt "gấp lên" và "chia đều":}\\
  - \textit{"Mỗi bạn có $5$ cái kẹo, có $4$ bạn"} $\implies$ phép nhân: $5 \times 4 = 20$.\\
  - \textit{"Có $20$ cái kẹo chia đều cho $4$ bạn"} $\implies$ phép chia: $20 : 4 = 5$.
  \item Trong phép nhân $a \times b$, nhớ viết thừa số chỉ số lượng mỗi phần ở phía trước.
\end{itemize}
}
\end{tcolorbox}
\end{minipage}


\end{document}
